\subsection{Subgroupoids}

DEFINITION 5. --- Let C be a category. A subcategory of C is a subquiver D of C satisfying the following conditions:

\begin{enumerate}
\def\labelenumi{(\roman{enumi})}
\item
  For every vertex a of D, \(e_a\) is an arrow of D;
\item
  For every composable pair \((f,g)\) of arrows of C belonging to D, the product fg is an arrow of D;
\end{enumerate}

Let C be a category and let D be a subcategory of C; endowed with the composition law induced from that of C by restriction to subsets, D is a category.

Every full subquiver of C is a subcategory of C. The intersection of a family of subcategories of C is a subcategory of C.

DEFINITION 6. --- Let G be a groupoid. A subgroupoid of G is a subcategory H of G such that the inverse of every arrow of G belonging to H is an arrow of H.

Let G be a groupoid. Every subgroupoid of G is a groupoid. Every full subquiver of G is a subgroupoid of G. The intersection of a family of subgroupoids of G is a subgroupoid of G.

Let H be a subquiver of G. The intersection of the subgroupoids of G of which H is a subquiver is called the subgroupoid of G generated by H. The set of its vertices is equal to that of H and its orbits are the connected components of H.

Examples. --- 1) Let X be a set; denote by X × X the groupoid of pairs of X (II, p. 163, example 2). Let R be an equivalence relation on X. The subquiver of the groupoid X × X whose set of vertices is X and whose set of arrows is the graph of the equivalence relation R is a subgroupoid of X × X. Its orbits are the equivalence classes of the relation R.

Conversely, every subgroupoid of X × X whose set of vertices is X is of this form.

Let X and S be sets and let \(p: X \rightarrow S\) be a map. The map p defines an equivalence relation on X (E, II, p. 41). We denote by \(X \times_{S} X\) the subgroupoid of \(X \times X\) defined by this equivalence relation and call it the groupoid of pairs of \((X, p)\) .

\begin{enumerate}
\def\labelenumi{\arabic{enumi})}
\setcounter{enumi}{1}
\item
  Let G and G' be groupoids, and let \(\varphi\) and \(\psi\) be morphisms of groupoids from G to G'. The equalizer of \(\varphi\) and \(\psi\) (cf. II, p. 153) is a subgroupoid of G.
\item
  Let X be a set and \(\Gamma\) a group. Denote by \(X \times \Gamma \times X\) the quiver whose set of vertices is X, whose set of arrows is \(X \times \Gamma \times X\) , the source and target maps being respectively \(\mathrm{pr}_1\) and \(\mathrm{pr}_3\) . Endowed with the composition law \((x, \gamma, x') \cdot (x', \gamma', x'') = (x, \gamma \gamma', x'')\) , this is a groupoid. It is moreover the groupoid deduced by transport of structure from the groupoid \((\mathrm{X}\times\mathrm{X})\times\Gamma\) , the product of the groupoid of pairs \(X\times X\) and the groupoid associated with the group \(\Gamma\) (II, p. 163, example 1), by means of the bijection \((x,x',\gamma)\mapsto(x,\gamma,x')\) from \(X\times X\times\Gamma\) to \(X\times\Gamma\times X\) .
\end{enumerate}

Let \(m\colon X \times \Gamma \to X\) be an operation (on the right) of the group \(\Gamma\) on \(X\) . The graph \(F\) of the map \(m\) , that is, the set of triples \((x, \gamma, x') \in X \times \Gamma \times X\) such that \(x' = x\gamma\) , is the set of arrows of a unique subgroupoid \(G\) of \(X \times \Gamma \times X\) . The orbits of this groupoid are the orbits of \(\Gamma\) on \(X\) ; if \(x \in X\) , the isotropy group of \(G\) at \(x\) is the set of \((x, \gamma, x)\) , where \(\gamma\) ranges over the stabilizer of \(x\) in \(\Gamma\) .

Conversely, any subgroupoid G of the groupoid \(X \times \Gamma \times X\) such that the map \((\mathrm{pr}_{1}, \mathrm{pr}_{2})\) from \(\mathrm{Fl}(\mathrm{G})\) to \(X \times \Gamma\) is a bijection is of this form.

\begin{enumerate}
\def\labelenumi{\arabic{enumi})}
\setcounter{enumi}{3}
\tightlist
\item
  Let G be a groupoid, let X be a set, and let \(\varphi\colon\mathrm{X}\to\mathrm{Som}(\mathrm{G})\) be a map. We define a quiver \(\mathrm{G}'\) as follows. The set \(\mathrm{Som}(\mathrm{G}')\) is the set X; for every pair \((x,y)\) of elements of X, the set \(\mathrm{Fl}_{x,y}(\mathrm{G}')\) is the set of triples \((x,f,y)\in\mathrm{X}\times\mathrm{Fl}(\mathrm{G})\times\mathrm{X}\) , where \(f\) is an element of \(\mathrm{Fl}_{\varphi(x),\varphi(y)}(\mathrm{G})\) . Let \(x,y,z\) be elements of X, \(f\in\mathrm{Fl}_{\varphi(x),\varphi(y)}(\mathrm{G})\) and \(g\in\mathrm{Fl}_{\varphi(y),\varphi(z)}(\mathrm{G})\) be arrows of G, and set \((x,f,y)\cdot(y,g,z)=(x,fg,z)\) . This defines a composition law in the quiver \(\mathrm{G}'\), making it a groupoid. It is called the inverse image groupoid of the groupoid G by the map \(\varphi\) and we denote it \(\varphi^{*}(\mathrm{G})\) .
\end{enumerate}

The pair \((\varphi,\psi)\) , where \(\psi\colon\mathrm{Fl}(\varphi^{*}(\mathrm{G}))\to\mathrm{Fl}(\mathrm{G})\) is the map defined by \((x,f,y)\mapsto f\) , is a morphism of groupoids from \(\varphi^{*}(\mathrm{G})\) to G, called the canonical morphism.

\begin{enumerate}
\def\labelenumi{\arabic{enumi})}
\setcounter{enumi}{4}
\tightlist
\item
  Let \(\varphi\colon\mathrm{G}\to\mathrm{G}^{\prime}\) be a morphism of groupoids. Let H be the subquiver of G with vertex set \(\mathrm{Som}(\mathrm{G})\) and whose set of arrows consists of the \(f\in\mathrm{Fl}(\mathrm{G})\) such that \(\varphi(f)\) is an identity element of \(\mathrm{G}^{\prime}\) . For all \(a\in\mathrm{Som}(\mathrm{G})\) , \(e_{a}\in\mathrm{H}_{a,a}\) . For all \(f\in\mathrm{Fl}(\mathrm{H})\) , \(\varphi(f^{-1})=\varphi(f)^{-1}\) ; therefore \(f^{-1}\in\mathrm{Fl}(\mathrm{H})\) . Moreover, if \(f\) and \(g\) are composable arrows of H, then \(\varphi(fg)=\varphi(f)\varphi(g)=e_{t(\varphi(f))}e_{o(\varphi(g))}=e_{t(\varphi(f))}\) , therefore \(fg\) is an arrow of H. This proves that H is a subgroupoid of G. It is called the kernel of \(\varphi\) and we denote it \(\mathrm{Ker}(\varphi)\) .
\end{enumerate}

\pandocbounded{\includegraphics[keepaspectratio]{images/part002_cf856079.jpg}}

\begin{enumerate}
\def\labelenumi{\arabic{enumi})}
\setcounter{enumi}{5}
\tightlist
\item
  Let \(\varphi\colon\mathrm{G}\to\mathrm{G}^{\prime}\) be a morphism of groupoids. The quiver \(\varphi(\mathrm{G})\) whose set of vertices is \(\varphi(\mathrm{Som}(\mathrm{G}))\) and whose set of arrows is \(\varphi(\mathrm{Fl}(\mathrm{G}))\) is in general not a subgroupoid of \(\mathrm{G}^{\prime}\) .
\end{enumerate}

This is nevertheless the case if the map Som( \(\varphi\) ) is injective (II, p. 225, exerc. 5).

\subsection{Operations of a groupoid}

Let G be a groupoid; denote by S the set of its vertices. Let X be a set and let \(p\colon \mathrm{X}\to \mathrm{S}\) be a map. A (right) operation \(\varphi\) of the groupoid G on the set X, relative to \(p\) , is a morphism of groupoids \(\varphi\) from G to the groupoid \(\mathcal{B}(\mathrm{X},p)\) of the permutations of the set X relative to \(p\) (II, p. 163) which induces the identity on the set of vertices. In other words, \(\varphi\) is the data, for every pair \((a,b)\) of points of S and for every arrow \(g\in \mathrm{G}\) joining \(a\) to \(b\), of a map \(\varphi (g)\colon \mathrm{X}_a\to \mathrm{X}_b\) , such that for every pair \((f,g)\) of composable arrows of G,

\[
\varphi (f g) = \varphi (g) \circ \varphi (f),
\]

and such that, for every \(a \in S\) , \(\varphi(e_a)\) is the identity of \(X_a\) .

Let \(\varphi\) and \(\varphi'\) be operations of the groupoid G on a set X, relative to a map \(p\colon X\to S\) . For \(\varphi=\varphi'\) , it suffices that there exist a subquiver H of G, generating G, such that \(\varphi(f)=\varphi'(f)\) for all \(f\in\mathrm{Fl}(\mathrm{H})\) . Indeed, the set of arrows \(f\) of G such that \(\varphi(f)=\varphi'(f)\) is the set of arrows of a subgroupoid of G, with vertex set S.

Let G be a groupoid with vertex set S, let X be a set, let \(p \colon \mathrm{X} \to \mathrm{S}\) be a map and let \(\varphi\) be an operation of G on X relative to \(p\) . Let \(\mathrm{G}_{\varphi}\) be the subquiver of \(p^{*}(\mathrm{G})\) whose set of vertices is X and whose set of arrows is the set of triples \((x, f, y) \in \mathrm{X} \times \mathrm{Fl}(\mathrm{G}) \times \mathrm{X}\) such that \(\varphi(f)(x) = y\) . It is a subgroupoid of \(p^{*}(\mathrm{G})\) .

Conversely, let \(\Gamma\) be a subgroupoid of \(p^{*}(\mathrm{G})\) ; suppose that the map \((x,f,y)\mapsto (x,f)\) from \(\mathrm{Fl}(\Gamma)\) to \(\mathrm{X}\times \mathrm{Fl}(\mathrm{G})\) is injective and that its image is the set of pairs \((x,f)\) such that \(p(x) = o(f)\) . There then exists a unique operation \(\varphi\) of the groupoid G on X such that one has \(\Gamma = \mathrm{G}_{\varphi}\) .

The orbits of the groupoid \(G_{\varphi}\) are called the orbits of the operation of G on X. By definition, they are the equivalence classes of the relation on X given by \(R\{x,y\}\) if and only if there exists \(f \in \mathrm{Fl}(G)\) such that \(\varphi(f)(x) = y\) .

Examples. --- 1) Let G be a group, let \(\mathcal{G}\) be the groupoid associated with G (II, p. 163, example 1). If X is a set, a (right) action of the groupoid \(\mathcal{G}\) on X is none other than a (right) action of the group G on X. Moreover, the orbits of the action of \(\mathcal{G}\) coincide with those of the group G.

\begin{enumerate}
\def\labelenumi{\arabic{enumi})}
\setcounter{enumi}{1}
\tightlist
\item
  Let \(\mathrm{G} = (\mathrm{S},\mathrm{F},o,t)\) be a groupoid. There exists a unique operation of \(\mathrm{G}\) on the set \(\mathrm{S}\) , relative to \(\mathrm{Id}_{\mathrm{S}}\) . It is given by \(\varphi(f)(a) = b\) if \(f \in \mathrm{G}_{a,b}\) . The orbits for this operation are the orbits of \(\mathrm{G}\) in the sense defined on p. 162.
\end{enumerate}

Let G be a groupoid, let S be the set of its vertices, let X be a set and let \(p: X \to S\) be a map. Let \(\varphi\) be an operation of G on X relative to p. One says that the groupoid G operates without monodromy on X if, for every point \(a \in S\) and every element \(f \in G_{a}\) , \(\varphi(f)\) is the identity of \(X_{a}\) . If the groupoid G is transitive, it suffices that this be so for one point of S.

Let \(a\) be a point of S and suppose that one has \(\varphi(f) = \mathrm{Id}_{\mathrm{X}_a}\) for all \(f \in \mathrm{G}_a\) . Let \(b\) be a point of S belonging to the orbit of \(a\) and let \(g\) be an arrow of G connecting \(a\) to \(b\) . For all \(f \in \mathrm{G}_b\) , \(\mathrm{Int}(g)(f) = gfg^{-1}\) is an element of \(\mathrm{G}_a\) ; one therefore has \(\mathrm{Id}_{\mathrm{X}_a} = \varphi(gfg^{-1}) = \varphi(g)\varphi(f)\varphi(g)^{-1}\) , so that \(\varphi(f) = \mathrm{Id}_{\mathrm{X}_b}\) . Let \(b\) be a point of S belonging to the orbit of \(a\) and let \(g\) , \(g'\) be arrows of G connecting \(a\) to \(b\) ; then \(g'g^{-1}\) is a loop at \(a\) , whence \(\varphi(g'g^{-1}) = \mathrm{Id}_{\mathrm{X}_a}\) and \(\varphi(g) = \varphi(g')\) .

\subsection{Distinguished subgroupoids; quotients of groupoids}

DEFINITION 7. --- Let G be a groupoid. One says that a subgroupoid H of G is distinguished if Som(H) = Som(G) and if, for every pair (a, b) of vertices of H and every arrow \(f \in \mathrm{G}_{a,b}\) , one has \(Int(f)(H_b) = H_a\) .

Let G be a groupoid and let H be a subgroupoid of G whose set of vertices is equal to Som(G). To verify that H is distinguished, it suffices to show that \(fH_{b}f^{-1} \subset H_{a}\) for all \(a \in \text{Som}(G)\) , all \(b \in \text{Som}(\mathbf{G})\) and every \(f \in \mathbf{G}_{a,b}\) . Indeed, if this is so, we also have \(f^{-1}\mathrm{H}_a f \subset \mathrm{H}_b\) for every arrow \(f \in \mathbf{G}_{b,a}\) , that is, \(\mathrm{H}_a \subset f\mathrm{H}_b f^{-1}\) , and, consequently, \(f\mathrm{H}_b f^{-1} = \mathrm{H}_a\) .

Let G be a groupoid and let H be a distinguished subgroupoid of G. For every vertex a of G, the subgroup \(H_{a}\) of \(G_{a}\) is a distinguished subgroup of \(G_{a}\) .

Let \(\varphi\colon\mathrm{G}\to\mathrm{G}^{\prime}\) be a morphism of groupoids. The kernel of \(\varphi\) (II, p. 166, example 5) is a distinguished subgroupoid of G. Indeed, let \(f\) be an arrow in G connecting \(a\) to \(b\) and let \(g\in\mathrm{Ker}(\varphi)_{b}\) ; we then have \(\varphi(fgf^{-1})=\varphi(f)\varphi(g)\varphi(f^{-1})=\varphi(f)e_{\varphi(b)}\varphi(f)^{-1}=e_{\varphi(a)}\) , whence the inclusion \(f\mathrm{Ker}(\varphi)_{b}f^{-1}\subset\mathrm{Ker}(\varphi)_{a}\) .

Let G be a groupoid. The groupoid G and the subgroupoid of G whose set of arrows is the set of identity elements of G are distinguished subgroupoids of G. The intersection of a family of distinguished subgroupoids of G is a distinguished subgroupoid of G. In particular, for every subset F ⊂ Fl(G), there exists a smallest distinguished subgroupoid of G whose set of arrows contains F. It is called the distinguished subgroupoid generated by F.

Let H be a distinguished subgroupoid of G. Let R be the equivalence relation in \(\mathrm{Fl}(\mathrm{G})\) defined by \(R\{f,g\}\) if and only if there exist arrows x and y in \(\mathrm{Fl}(\mathrm{H})\) such that f = xgy. If f and g are arrows of G equivalent modulo R, their origins (resp. their targets) belong to the same orbit of H. Set \(F' = \mathrm{Fl}(\mathrm{G}) / \mathcal{R}\) and denote by \(o'\) and \(t'\) the maps from \(F'\) to \(\mathrm{Orb}(\mathrm{H})\) deduced from o and t by passing to quotients. Denote by G/H the quiver \((\mathrm{Orb}(\mathrm{H}), \mathrm{F}', o', t')\) .

Lemma. --- Given composable arrows u and v of G/H, one can choose representatives f and g of u and v in Fl(G) which are composable. The class modulo R of the product fg does not depend on the choice of f and g.

Let \(f\) and \(g\) be arbitrary representatives of \(u\) and \(v\) in \(\mathrm{Fl(G)}\) . The target of \(f\) and the origin of \(g\) belong to the same orbit of H, hence can be connected by an arrow \(x\) of H (II, p. 162, prop. 1). The arrows \(fx\) and \(g\) are then representatives of \(u\) and \(v\) in F that are composable in G. This proves the first assertion.

Now let \(f\) , \(g\) on the one hand, and \(f'\) , \(g'\) on the other hand, be representatives of \(u\) and \(v\) which are composable in G. By hypothesis, there exist arrows \(x\) , \(y\) , \(z\) , \(t\) in H such that one has \(f' = xfy\) and \(g' = zgt\) . One then has \(yz \in \mathrm{H}_{t(f)}\) and \(f'g' = x f y z g t = x f(yz)f^{-1}f g t\) . Since H is a distinguished subgroupoid of G, the arrow \(f(yz)f^{-1}\) is an arrow of H. The same is therefore true of the arrow \(x f(yz)f^{-1}\) , which proves that \(f'g'\) and \(fg\) are equivalent modulo \(\mathcal{R}\) .

By virtue of this lemma, there exists a unique composition law \(m\) in the quiver G/H such that, for every pair \((u,v)\) of composable arrows, \(m(u,v)\) is the class modulo \(\mathcal{R}\) of the product \(fg\) , for every pair \((f,g)\) of composable arrows of G such that \(u\) is the class of \(f\) and \(v\) that of \(g\) . Endowed with this composition law, G/H is a groupoid. Let \(p_1\colon \operatorname{Som}(\mathrm{G}) \to \operatorname{Som}(\mathrm{G}/\mathrm{H})\) and \(p_2\colon \operatorname{Fl}(\mathrm{G}) \to \operatorname{Fl}(\mathrm{G}/\mathrm{H})\) be the canonical surjections. The pair \(p = (p_1,p_2)\) is a groupoid morphism from G to G/H whose kernel is the distinguished subgroupoid H.

DEFINITION 8. --- We say that G/H is the quotient groupoid of G by H and that p: G → G/H is the canonical morphism.

Remarks. --- 1) The map Som(p) defines, by passing to quotients, a bijection from the set of orbits of G to the set of orbits of G/H. In particular, G is transitive if and only if G/H is.

\begin{enumerate}
\def\labelenumi{\arabic{enumi})}
\setcounter{enumi}{1}
\tightlist
\item
  Let \(a\) and \(b\) be vertices of G. The map from \(\mathrm{G}_{a,b}\) to \((\mathrm{G} / \mathrm{H})_{p(a),p(b)}\) deduced from \(p\) is surjective. Indeed, let \(u\) be an arrow from G/H connecting \(p(a)\) to \(p(b)\) ; it is the class modulo \(\mathcal{R}\) of an arrow \(f\) of G. The origin \(o(f)\) of \(f\) belongs to the orbit of \(a\) in H; there therefore exists an arrow \(x\) in H connecting \(a\) to \(o(f)\) . Similarly, there exists an arrow \(y\) in H that connects \(t(f)\) to \(b\) . Then, \(f' = xfy\) is an element of \(\mathrm{G}_{a,b}\) whose class modulo \(\mathcal{R}\) is \(u\) .
\end{enumerate}

PROPOSITION 2. --- Let a be a vertex of G. The group homomorphism \(p_{a}\colon\mathrm{G}_{a}\to(\mathrm{G}/\mathrm{H})_{p(a)}\) deduced from p by passing to isotropy groups is surjective; its kernel is \(\mathrm{H}_{a}\) .

The homomorphism \(p_a\) is surjective by virtue of the preceding remark. Its kernel contains \(\mathrm{H}_a\) by construction of the groupoid \(\mathrm{G / H}\) . Let \(f\) be an element of \(\mathrm{G}_a\) such that \(p_a(f) = e_{p(a)}\) . This means that \(f\) and \(e_a\) are equivalent modulo \(\mathcal{R}\) ; there then exist arrows \(x\) and \(y\) of H such that \(f = xe_a y\) . Necessarily, \(x\) and \(y\) belong to \(\mathrm{H}_a\) , whence \(f \in \mathrm{H}_a\) .

PROPOSITION 3. --- Let G be a groupoid, let H be a distinguished subgroupoid of G, and let p: G → G/H be the canonical morphism. Let \(\varphi\colon G\to G'\) be a morphism of groupoids such that \(H\subset\operatorname{Ker}(\varphi)\) . There exists a unique groupoid morphism \(\overline{\varphi}\colon G/H\to G'\) such that \(\overline{\varphi}\circ p=\varphi\) .

We say that \(\overline{\varphi}\) is the groupoid morphism induced by \(\varphi\) by passing to the quotient.

The uniqueness of such a morphism is evident, since the maps \(\operatorname{Som}(p)\) and \(\operatorname{Fl}(p)\) are surjective.

Let \(a\) and \(b\) be vertices of G. If they are in the same orbit of H, there exists an arrow \(f\) connecting \(a\) to \(b\) in H and one has \(\varphi(f) = e_{\varphi(a)}\) . In particular, \(\varphi(a) = \varphi(b)\) . Consequently, the map \(\operatorname{Som}(\varphi)\) defines, by passing to the quotient, a map \(\overline{\varphi}_1\colon \operatorname{Orb}(\mathrm{H}) \to \operatorname{Som}(\mathrm{G}')\) . Let \(f\) and \(g\) be arrows in G. If \(f\) and \(g\) are equivalent modulo \(\mathcal{R}\) , there exist arrows \(x\) and \(y\) in H such that \(f = xgy\) . Consequently, \(\varphi(f) = \varphi(x)\varphi(g)\varphi(y) = \varphi(g)\) since \(\varphi(x)\) and \(\varphi(y)\) are identity elements. The map \(\operatorname{Fl}(\varphi)\) therefore defines, by passing to the quotient, a map \(\overline{\varphi}_2\colon \operatorname{Fl}(\mathrm{G}) / \mathcal{R} \to \operatorname{Fl}(\mathrm{G}')\) .

Let \(f\) and \(g\) be arrows of \(G\) which are composable; denote by \(u\) and \(v\) their classes in \(\mathrm{Fl}(\mathrm{G} / \mathrm{H})\) . We have \(\overline{\varphi}_2(uv) = \varphi(fg) = \varphi(f)\varphi(g) = \overline{\varphi}_2(u)\overline{\varphi}_2(v)\) .

The pair \(\overline{\varphi} = (\overline{\varphi}_1, \overline{\varphi}_2)\) is a morphism of groupoids from G/H to G', and one has \(\overline{\varphi} \circ p = \varphi\) .

\subsection{Groupoid of path classes of a graph}

Let G be a graph. Denote by \(\Omega_{\mathrm{G}}\) the category of paths of the quiver underlying G (II, p. 160, example 1). Consider the finest equivalence relation \(\mathcal{R}\) in \(\mathrm{Ch(G)}\) such that:

\begin{enumerate}
\def\labelenumi{(\roman{enumi})}
\item
  For every pair \((a,b)\) of vertices of G and every arrow f in G connecting a to b, the paths \((a,f,b,\overline{f},a)\) and \(e_{a}\) are equivalent modulo R;
\item
  If \((c,d)\) and \((c',d')\) are pairs of juxtaposable paths in G such that \(R\{c,c'\}\) and \(R\{d,d'\}\) are equivalent, then the paths \(c*d\) and \(c'*d'\) are equivalent modulo R.
\end{enumerate}

Two paths equivalent modulo \(\mathcal{R}\) have the same origin and the same target. The maps \(o\) and \(t\) of \(\mathrm{Ch(G)}\) to \(\mathrm{Som(G)}\) define, by passing to the quotient, maps \(o'\) and \(t'\) from \(\mathrm{Ch(G)/\mathcal{R}}\) to \(\mathrm{Som(G)}\) . Denote by \(\varpi_{G}\) the quiver (Som(G), Ch(G)/R, \(o', t'\) ). The pair \(\varphi\) formed by the identity of Som(G) and the canonical projection of Ch(G) onto Ch(G)/R is a quiver morphism from \(\Omega_{G}\) to \(\varpi_{G}\) . The juxtaposition of paths in Ch(G) defines, by passing to quotients, a composition law in \(\varpi_{G}\) .

PROPOSITION 4. --- Equipped with this composition law, \(\varpi_{\mathrm{G}}\) is a groupoid.

By construction, we have the relation \(\varphi(cc') = \varphi(c)\varphi(c')\) , for every pair of paths \((c,c')\) juxtaposable in G. Every arrow of \(\varpi_{\mathrm{G}}\) is of the form \(\varphi(c)\) , where \(c\) is a path in G. This implies that the composition law of \(\varpi_{\mathrm{G}}\) is associative and that \(\varphi(e_a)\) is a neutral element at \(a\) for every vertex \(a\) of \(\varpi_{\mathrm{G}}\) . It remains to show that every arrow of \(\varpi_{\mathrm{G}}\) is invertible. Let \(c\) be a path in G and let us prove by induction on the length of \(c\) that one has \(\varphi(c)\varphi(\overline{c}) = \varphi(e_{o(c)})\) . This equality is true if \(c\) has length 0. If \(c\) has length \(n \geqslant 1\) , one can write \(c = c_1c_2\) , with \(c_1\) of length 1 and \(c_2\) of length \(n - 1\) . Then, \(\overline{c} = \overline{c_2}\overline{c_1}\) and one has

\[
\begin{array}{c} \varphi (c) \varphi (\overline {{c}}) = \varphi (c _ {1}) \varphi (c _ {2}) \varphi (\overline {{c _ {2}}}) \varphi (\overline {{c _ {1}}}) = \varphi (c _ {1}) \varphi (e _ {o (c _ {2})}) \varphi (\overline {{c _ {1}}}) = \varphi (c _ {1}) \varphi (\overline {{c _ {1}}}) \\ = \varphi (c _ {1} \overline {{c _ {1}}}) = \varphi (e _ {o (c _ {1})}) \end{array}
\]

by definition of the relation \(\mathcal{R}\) .

Applying this equality to the path \(\overline{c}\) , we see that \(\varphi(\overline{c})\varphi(c)=\varphi(e_{t(c)})\) . Thus, \(\varphi(c)\) is invertible, with inverse \(\varphi(\overline{c})\) .

The equivalence classes of the relation \(\mathcal{R}\) are called the path classes in the graph G; the groupoid \(\varpi_{\mathrm{G}}\) is called the groupoid of path classes of the graph G. It has the same vertex set as G and its orbits are the connected components of the graph G.

Denote by \(v\) the map which, to an arrow \(f\) of \(G\) , associates the class of the path \((o(f), f, t(f))\) of \(G\) . The pair \(j = (\mathrm{Id}_{\mathrm{Som}(G)}, v)\) is a morphism, called canonical, of quivers from \(G\) to \(\varpi_G\) . Its image generates \(\varpi_G\) ; for every arrow \(f\) of \(G\) , we have \(j(\overline{f}) = j(f)^{-1}\) .

Let G and \(G'\) be graphs, let \(\varphi\colon G \to G'\) be a graph morphism. Denote by \(j\colon G \to \varpi_{G}\) and \(j'\colon G' \to \varpi_{G'}\) the canonical morphisms. If \(c = (a_0, f_1, a_1, \ldots, a_n)\) is a path in G, the class of the path \(\varphi(c) = (\varphi(a_0), \varphi(f_1), \varphi(a_1), \ldots, \varphi(a_n))\) depends only on the class of c. One thus defines, by passing to equivalence classes, a quiver morphism \(\varpi(\varphi)\colon\varpi_{\mathrm{G}}\to\varpi_{\mathrm{G}^{\prime}}\) such that \(\varpi(\varphi)\circ j=j^{\prime}\circ\varphi\) . This is a morphism of groupoids.

PROPOSITION 5. — Let G be a graph; denote by j the canonical quiver morphism from G to \(\varpi_{\mathrm{G}}\) . Let \(\varphi\) be a quiver morphism from G into a groupoid \(\mathrm{G}'\) such that \(\varphi(\overline{f}) = \varphi(f)^{-1}\) for every arrow \(f\) of G. Then there exists a unique morphism of groupoids \(\varphi'\) from \(\varpi_{\mathrm{G}}\) to \(\mathrm{G}'\) such that \(\varphi' \circ j = \varphi\) .

One defines a map \(u\colon \mathrm{Ch}(\mathrm{G})\to \mathrm{Fl}(\mathrm{G}^{\prime})\) by setting, for every path \(c = (a_0,f_1,a_1,\ldots ,f_n,a_n)\) in G, \(u(c) = e_{a_0}\varphi (f_1)\dots \varphi (f_n)\) . For every pair \((c,c')\) of juxtaposable paths in G, we have \(u(cc^{\prime}) = u(c)u(c^{\prime})\) . The map \(u\) is compatible with the equivalence relation \(\mathcal{R}\) defined above, whence, by passing to the quotient, a map \(u'\colon \mathrm{Ch}(\mathrm{G}) / \mathcal{R}\to \mathrm{Fl}(\mathrm{G}')\) . Let us then set \(\varphi ' = (\mathrm{Som}(\varphi),u')\) . This is a groupoid morphism from \(\varpi_{\mathrm{G}}\) to \(\mathrm{G}'\) such that \(\varphi '\circ j = \varphi\) .

Let \(\psi\) be a groupoid morphism from \(\varpi_{\mathrm{G}}\) to \(\mathbf{G}'\) such that \(\psi \circ j = \varphi\) . The equalizer of \(\varphi'\) and \(\psi\) is a subgroupoid of \(\varpi_{\mathrm{G}}\) which contains \(j(\mathrm{G})\) . It is therefore equal to \(\varpi_{\mathrm{G}}\) since \(\varpi_{\mathrm{G}}\) is generated by \(j(\mathrm{G})\) and one has \(\psi = \varphi'\) .

COROLLARY 1. --- In a graph, every path class contains a unique path without backtracking.

Let G be a graph. Denote by \(j\colon G\to\varpi_{G}\) the canonical quiver morphism.

The existence, for every path \(c\) of G, of an equivalent path without backtracking is immediate by induction on the length of \(c\) (cf. II, p. 157).

Let A be an orientation of G and let G' be the groupoid associated with the free group F(A) constructed on A (II, p. 163, example 1). Let ψ be the quiver morphism from G to G' such that, for every \(f \in A\) , \(\psi(f)\) is the element f of F(A) and \(\psi(\overline{f})\) is the element \(f^{-1}\) of F(A). Let ψ' be the unique groupoid morphism from \(\varpi_{G}\) to G' such that ψ' ∘ j = ψ.

Let \(c\) and \(c'\) be paths without backtracking equivalent modulo \(\mathcal{R}\) . They have the same source and the same target. To prove that they are equal, it suffices to prove that the sequences \((f_1, \ldots, f_n)\) and \((g_1, \ldots, g_m)\) of their arrows are equal. The image under \(\psi'\) of the common class of \(c\) and \(c'\) is equal to \(\psi(f_1) \ldots \psi(f_n)\) and to \(\psi(g_1) \ldots \psi(g_m)\) . Now the terms of the sequences \((\psi(f_1), \ldots, \psi(f_n))\) and \((\psi(g_1), \ldots, \psi(g_m))\) belong to the subset \(\mathrm{A} \cup \mathrm{A}^{-1}\) of \(\mathrm{F(A)}\) and two consecutive elements of these sequences are not inverses of one another. By A, I, p. 84, prop. 7, these two sequences are equal. It follows that the sequences \((f_{1},\ldots,f_{n})\) and \((g_{1},\ldots,g_{m})\) are equal, and hence that \(c=c'\) , whence uniqueness.

COROLLARY 2. --- The canonical morphism from G to \(\varpi_{\mathrm{G}}\) is injective.

This follows immediately from Corollary 1.

COROLLARY 3. --- Let G' be a subgraph of G. The morphism from \(\varpi_{\mathrm{G}'}\) to \(\varpi_{\mathrm{G}}\) deduced from the injection of G' into G is injective.

Denote by \(i\) the inclusion morphism of \(\mathbf{G}'\) into \(\mathbf{G}\) and \(\varpi(i)\) the morphism from \(\varpi_{\mathbf{G}'}\) to \(\varpi_{\mathbf{G}}\) which is deduced from it. The maps \(\operatorname{Som}(i)\) and \(\operatorname{Som}(\varpi(i))\) coincide. Moreover, if \(c\) is a path with no backtracking in \(\mathbf{G}'\) , the path \(i(c)\) is a path with no backtracking in \(\mathbf{G}\) . The assertion therefore follows from Corollary 1.

COROLLARY 4. --- A nonempty graph G is a tree if and only if the groupoid \(\varpi_{\mathrm{G}}\) is simply transitive.

Let \(a\) be a point of G. By Corollary 1, the following properties are equivalent:

\begin{enumerate}
\def\labelenumi{(\roman{enumi})}
\item
  The isotropy group of \(\varpi_{\mathrm{G}}\) at \(a\) is reduced to the identity element;
\item
  The only loop of G originating from \(a\) which is without backtracking is the constant loop based at \(a\) .
\end{enumerate}

The groupoid \(\varpi_{G}\) is simply transitive if and only if it is transitive and has property (i) for every a. The graph G is a tree if and only if it is connected and has property (ii) for every point a. Since the orbits of \(\varpi_{G}\) are identified with the connected components of G, the corollary follows.

\subsection{Free groupoids}

DEFINITION 9. --- Let C be a quiver. The free groupoid constructed on C, denoted \(\mathrm{Grp}(\mathrm{C})\), is the groupoid \(\varpi_{\widetilde{C}}\) of path classes of the graph \(\widetilde{C}\) associated with C.

The set of vertices of Grp(C) is equal to that of C; the orbits of Grp(C) are the connected components of C.

Let C be a nonempty quiver. By Corollary 4, p. 174, the graph \(\widetilde{\mathrm{C}}\) associated with C is a tree if and only if the free groupoid \(\mathrm{Grp}(\mathrm{C})\) constructed on C is simply transitive.

The composite of the canonical quiver morphisms \(i\colon \mathrm{C}\to \widetilde{\mathrm{C}}\) and \(j\colon \widetilde{\mathrm{C}}\to \varpi_{\widetilde{\mathrm{C}}}\) is a quiver morphism from C to \(\mathrm{Grp}(\mathrm{C})\) which we call the canonical morphism from C to \(\mathrm{Grp}(\mathrm{C})\) . Let us denote it \(\theta\) . For every vertex \(a\) of C, we have \(\theta (a) = a\) . If \(f\) is an arrow of C, \(\theta (f)\) is the class of the path \((o(f),(f,1),t(f))\) of \(\widetilde{\mathrm{C}}\) . The morphism \(\theta\) is injective (II, p. 174, cor. 2) and its image generates the groupoid \(\mathrm{Grp}(\mathrm{C})\) .

Let \(\varphi\colon\mathrm{C}\to\mathrm{C}^{\prime}\) be a quiver morphism. Denote by \(\theta_{\mathrm{C}}\) (resp. \(\theta_{\mathrm{C}^{\prime}}\) ) the canonical morphism from C to \(\mathrm{Grp}(\mathrm{C})\) (resp. from \(\mathrm{C}^{\prime}\) to \(\mathrm{Grp}(\mathrm{C}^{\prime})\) ). The morphism \(\varpi(\widetilde{\varphi})\) is the unique groupoid morphism \(\psi\) from \(\mathrm{Grp}(\mathrm{C})\) to \(\mathrm{Grp}(\mathrm{C}^{\prime})\) such that \(\psi\circ\theta_{\mathrm{C}}=\theta_{\mathrm{C}^{\prime}}\circ\varphi\) .

PROPOSITION 6. --- Let C be a quiver, let G be a groupoid, and let \(\varphi\) be a quiver morphism from C to G. There exists a unique groupoid morphism \(\varphi'\) from Grp(C) to G such that \(\varphi' \circ \theta_{\mathrm{C}} = \varphi\) .

Let \(\psi\) be the quiver morphism from \(\widetilde{\mathbf{C}}\) to G such that \(\psi(a) = \varphi(a)\) for every vertex \(a\) of C and \(\psi(f, \varepsilon) = \varphi(f)^{\varepsilon}\) for every arrow \(f\) of C and every \(\varepsilon \in \{-1, 1\}\) . We have \(\psi \circ i = \varphi\) . By Proposition 5, p. 173, there exists a morphism \(\varphi'\) from \(\mathrm{Grp}(\mathrm{C}) = \varpi_{\widetilde{\mathrm{C}}}\) to G such that \(\varphi' \circ j = \psi\) . One then has \(\varphi' \circ \theta_{\mathrm{C}} = \varphi' \circ j \circ i = \psi \circ i = \varphi\) . As in the proof of Prop. 5, the uniqueness of \(\varphi'\) results from the fact that \(\theta_{\mathrm{C}}(\mathrm{C})\) generates the groupoid \(\mathrm{Grp}(\mathrm{C})\) .

Under the hypotheses of Prop. 6, one sometimes says that \(\varphi'\) is the groupoid morphism from Grp(C) to G extending the quiver morphism \(\varphi\) .

Example. --- Let C be a quiver with a single vertex s and let A be the set of its arrows. The groupoid \(\mathrm{Grp}(\mathrm{C})\) is then the groupoid associated with the free group \(\mathrm{F}(\mathrm{A})\) built on A. Indeed, it has a single vertex s, and it follows immediately from Prop. 6 that the canonical map from A to \(\mathrm{Grp}(\mathrm{C})_{s}\) deduced from the canonical morphism from C to \(\mathrm{Grp}(\mathrm{C})\) satisfies the universal property of free groups (A, I, p. 85, prop. 8).

\subsection{Contraction of arrows of a groupoid}

Let G be a groupoid and let F be a subset of the set of arrows of G. Let H be the distinguished subgroupoid of G generated by F. The groupoid G/H is called the groupoid obtained from G by contracting the set of arrows F and is denoted G/F. If p denotes the canonical morphism from G to G/F, we have \(p(o(f)) = p(t(f))\) and \(p(f) = e_{p(o(f))}\) for every arrow \(f \in F\) .

This morphism satisfies the following universal property.

PROPOSITION 7. --- For any morphism of groupoids \(\varphi\colon G\to G'\) which maps every arrow belonging to F to an identity element of \(G'\) , there exists a unique morphism of groupoids \(\varphi'\colon G/F\to G'\) such that \(\varphi'\circ p=\varphi\) .

The kernel of \(\varphi\) is a distinguished subgroupoid of G, and the set of its arrows contains F. By definition, H is the smallest distinguished subgroupoid of G whose set of arrows contains F; it is therefore contained in the kernel of \(\varphi\) . The proposition then follows from II, p. 170, prop. 3.

Denote by \(\Gamma\) the subquiver of G with vertex set Som(G) and arrow set F. The orbits of H are identified with the connected components of the quiver \(\Gamma\) . By definition of the quotient of a groupoid by a distinguished subgroupoid, the vertex set of G/F is identified with the set of connected components of \(\Gamma\) ; in other words, it is the quotient set of Som(G) by the finest equivalence relation such that \(o(f)\) is equivalent to \(t(f)\) for every arrow \(f \in \mathrm{F}\) .

The map p induces, by passing to quotients, a bijection from the set of orbits of G to the set of orbits of G/F (II, p. 170, remark 1).

The purpose of this section is to compute the homomorphisms induced by p by passing to the isotropy groups, which, according to Prop. 2 (II, p. 170), amounts to computing the isotropy groups of the subgroupoid H.

PROPOSITION 8. --- Let δ be the canonical morphism from the free groupoid \(\mathrm{Grp}(\Gamma)\) into G extending the canonical injective morphism from Γ into G. Let a be a vertex of G and let A be the orbit of a in G. For every b ∈ A, let \(f_{ab}\) be an arrow of G connecting a to b. The isotropy group \(H_{a}\) is the distinguished subgroup of \(G_{a}\) generated by the elements \(\operatorname{Int}(f_{ab})(\delta(c))\) , where b ranges over the set A and c ranges over the elements of \(\mathrm{Grp}(\Gamma)_{b}\) .

If x and y are vertices of G belonging to the same orbit, with \(x \neq a\) ; we further fix an arrow \(f_{xy}\) of G that connects x to y. For every point x of G, let \(N_{x}\) be the distinguished subgroup of \(G_{x}\) generated by the elements \(\operatorname{Int}(f_{xy})(\delta(c))\) , where y ranges over the orbit of x in G and c ranges over the group \(\operatorname{Grp}(\Gamma)_{y}\) . For all \(f \in G_{xy}\) and every \(g \in N_{y}\) , we have \(fgf^{-1} \in N_{x}\) . Let \(x \in \operatorname{Som}(G)\) . We have \(\delta(\operatorname{Grp}(\Gamma)_{x}) \subset N_{x}\) , by definition of \(N_{x}\) . By definition of the groupoid H, \(\delta(f)\) is an arrow of H for every arrow f of \(\operatorname{Grp}(\Gamma)\) ; consequently, \(\delta(\operatorname{Grp}(\Gamma)_{x})\) is contained in \(H_{x}\) . Since H is a distinguished subgroupoid of G, we then have

\[
\operatorname{Int} (f _ {x y}) (\delta (c)) \in \operatorname{Int} (f _ {x y}) (\mathrm{N} _ {y}) \subset \operatorname{Int} (f _ {x y}) (\mathrm{H} _ {y}) \subset \mathrm{H} _ {x}
\]

for all \(y \in \operatorname{Som}(\mathbf{G})\) , then \(\mathrm{N}_x \subset \mathrm{H}_x\) .

Let \(x\) and \(y\) be vertices of \(\Gamma\) . Set \(\mathrm{N}_{x,y} = \varnothing\) if \(x\) and \(y\) do not belong to the same connected component of \(\Gamma\) . Otherwise, let \(c\) and \(c'\) be path classes connecting \(x\) to \(y\) in the graph \(\widetilde{\Gamma}\) associated with \(\Gamma\) . Let \(z\) be a vertex of G belonging to the orbit of \(y\) and let \(\ell\) be an element of \(\mathrm{Grp}(\Gamma)_z\) ; in the group \(\mathrm{G}_x\) one has the equality:

\[
\begin{array}{r l} & {\delta (c) f _ {y z} \delta (\ell) f _ {y z} ^ {- 1} \delta (c ^ {\prime}) ^ {- 1}} \\ & {\qquad = \delta (c) f _ {y z} f _ {x z} ^ {- 1} \big (f _ {x z} \delta (\ell) f _ {x z} ^ {- 1} \big) f _ {x z} f _ {y z} ^ {- 1} \delta (c) ^ {- 1} \delta (c (c ^ {\prime}) ^ {- 1}).} \end{array}
\]

The arrow \(\delta(c(c')^{-1})\) is the class of a loop at \(x\) in the quiver \(\Gamma\) , hence belongs to \(\mathrm{N}_x\) , just as the arrow \(f_{xz}\delta(\ell)f_{xz}^{-1}\) . Since \(\delta(c)f_{yz}f_{xz}^{-1}\) belongs to \(\mathrm{G}_x\) and \(\mathrm{N}_x\) is a distinguished subgroup of \(\mathrm{G}_x\) , it follows that \(\delta(c)f_{yz}\delta(\ell)f_{yz}^{-1}\delta(c')^{-1}\) belongs to \(\mathrm{N}_x\) . Consequently, \(\delta(c)\mathrm{N}_y\subset \mathrm{N}_x\delta(c')\) . By symmetry, we have \(\delta(c)\mathrm{N}_y = \mathrm{N}_x\delta(c')\) . This implies that this set does not depend on the choices of \(c\) and \(c'\) ; it is denoted \(\mathrm{N}_{x,y}\) .

Let N be the subquiver of G whose set of vertices is Som(G) and whose set of arrows from x to y is equal to \(N_{x,y}\) for every pair \((x,y)\) of vertices of G. It is a distinguished subgroupoid of G whose set of arrows contains F; consequently, H is a subgroupoid of N. In particular, \(H_{a} \subset N_{a}\) , whence the equality.

COROLLARY 1. --- Suppose further that, for every arrow f of F, we have \(o(f) = t(f)\) . Then, the origin and the target of every path in the quiver \(\Gamma\) coincide. Let a be a vertex of G and let A be its orbit in G. For every element c of F whose origin b belongs to A, fix an arrow \(f_{c}\) of G connecting a to b, and set \(\kappa(c) = \operatorname{Int}(f_{c})(\delta(c)) = f_{c}\delta(c)f_{c}^{-1}\) .

The group \(H_{a}\) is then the distinguished subgroup of \(G_{a}\) generated by the arrows \(\kappa(c)\) .

By prop. 8, the arrows \(\kappa(c)\) are contained in the group \(H_a\) . Let \(N_a\) be the smallest normal subgroup of \(G_a\) containing them. By this proposition, it suffices to prove that, for every vertex \(b\) of G belonging to A, every arrow \(f_{ab}\) connecting \(a\) to \(b\) in G and every element \(c\) of \(\mathrm{Grp}(\Gamma)_b\) , \(\operatorname{Int}(f_{ab})(\delta(c))\) belongs to \(N_a\) . Since the origin and the target of each element of F are equal, a loop at a vertex \(b\) in the graph \(\widetilde{\Gamma}\) is of the form \((b,(f_1,\varepsilon_1),b,\ldots,(f_n,\varepsilon_n),b)\), where for every \(i\in\{1,\ldots,n\}\), \(f_i\) is an element of F with origin \(b\) and \(\varepsilon_i\in\{-1,1\}\). We have

\[
\operatorname{Int} (f _ {a b}) (\delta (c)) = f _ {a b} f _ {c} ^ {- 1} \kappa (f _ {1}) ^ {\varepsilon_ {1}} \dots \kappa (f _ {n}) ^ {\varepsilon_ {n}} f _ {c} f _ {a b} ^ {- 1}.
\]

Since \(f_{ab}f_c^{-1} \in \mathrm{G}_a\) , it follows that \(\operatorname{Int}(f_{ab})(\delta(c))\) belongs to \(\mathrm{N}_a\) .

COROLLARY 2. --- If the graph associated with the quiver \(\Gamma\) is a forest, the homomorphism \(p_a\) from \(\mathrm{G}_a\) to \((\mathrm{G} / \mathrm{F})_{p(a)}\) is bijective, for every vertex a of G.

Indeed, under this hypothesis, we have \(\mathrm{Grp}(\Gamma)_{b}=\{e_{b}\}\) for every vertex b of G (corollary 1 of II, p. 173). It then follows from Proposition 8 that, for every vertex a of G, the group \(H_{a}\) is reduced to the identity element.

\subsection{Poincaré group of a graph}

Let G be a graph. Let a be a vertex of G; the isotropy group at a of the groupoid \(\mathrm{Grp}(G)\) is called the Poincaré group of G at a and is denoted \(\pi_{1}(G,a)\) . Let c be a path class in G, let a be its origin and b its target. The map \(\operatorname{Int}(c): \pi_{1}(G,b)\to\pi_{1}(G,a)\) defined by \(c'\mapsto cc'c^{-1}\) is a group isomorphism. Let \(\varphi:G\to H\) be a graph morphism. Denote by \(\theta_{G}:G\to\mathrm{Grp}(G)\) and \(\theta_{H}:H\to\mathrm{Grp}(H)\) the canonical morphisms. If \(\overline{\varphi}\) denotes the unique groupoid morphism \(\mathrm{Grp}(G) \to \mathrm{Grp}(H)\) such that \(\overline{\varphi}\circ\theta_{G}=\theta_{H}\circ\varphi\) , the group homomorphism \(\overline{\varphi}_{a}:\pi_{1}(G,a)\to\pi_{1}(H,\varphi(a))\) is denoted \(\pi_{1}(\varphi,a)\) .

Let G be a connected graph, let S be an orientation of G, and let A be a maximal oriented tree of G (II, p. 157, prop. 1). Given vertices \(a\) and \(b\) of G, there exists a unique path class \(\gamma_{a,b}\) in the graph \(\widetilde{\mathrm{A}}\) associated with A which connects \(a\) to \(b\) (II, p. 174, cor. 4 of prop. 5). If \(a\) , \(b\) and \(c\) are vertices of G, we have \(\gamma_{a,b}\gamma_{b,c} = \gamma_{a,c}\) , these two classes of paths being equal to the unique class of paths connecting \(a\) to \(c\) in \(\widetilde{\mathrm{A}}\) .

PROPOSITION 9. --- Let a be a vertex of G. There exists a unique homomorphism \(\lambda\) from the free group \(\mathrm{F}(\mathrm{S}-\mathrm{Fl}(\mathrm{A}))\) to \(\pi_1(G, a)\) such that

\[
\lambda (f) = \gamma_ {a, o (f)} \cdot f \cdot \gamma_ {t (f), a}\tag{1}
\]

for \(f \in \mathrm{S} - \mathrm{Fl}(\mathrm{A})\) . The homomorphism \(\lambda\) is a group isomorphism.

Let L denote the group \(\mathrm{F}(\mathrm{S}-\mathrm{Fl}(\mathrm{A}))\) . The existence and uniqueness of the homomorphism \(\lambda\colon L\to\pi_{1}(\mathrm{G},a)\) satisfying relations (1) follows from A, I, p. 85, prop. 8. Let L be the groupoid associated with the group L (II, p. 163, example 1); denote by s its unique vertex. There exists a unique groupoid morphism \(\mu\colon\mathrm{Grp}(\mathrm{G})\to\mathcal{L}\) such that \(\mu(\theta_{\mathrm{G}}(f))=f\) for every arrow \(f\in\mathrm{S}-\mathrm{Fl}(\mathrm{A})\) and \(\mu(\theta_{\mathrm{G}}(f))=e_{s}\) for every arrow \(f\in\mathrm{Fl}(\mathrm{A})\) (II, p. 173, prop. 5). For \(f\in\mathrm{S}-\mathrm{Fl}(\mathrm{A})\) , we have \(\mu(\theta_{\mathrm{G}}(f))=f\) ; the homomorphism \(\mu_{a}\circ\lambda\) is therefore the identity isomorphism of the group L (A, I, p. 85, prop. 8). It follows that \(\lambda\) is injective.

Let \(\varpi_{\mathrm{G}} / \mathrm{A}\) be the groupoid deduced from \(\varpi_{\mathrm{G}}\) by contraction of the arrows of A and let \(p\colon \varpi_{\mathrm{G}}\to \varpi_{\mathrm{G}} / \mathrm{A}\) be the canonical groupoid morphism. It is surjective, and the group homomorphism \(p_a\) deduced from \(p\) by passing to isotropy groups is an isomorphism (II, p. 178, cor. 2 of prop. 8). Since the graph G is assumed connected and A is a maximal tree thereof, the groupoid \(\varpi_{\mathrm{G}} / \mathrm{A}\) has a unique vertex (II, p. 176). Moreover, \(\varpi_{\mathrm{G}}\) is generated by \(\theta_{\mathrm{G}}(\mathrm{G})\) ; therefore, \(\varpi_{\mathrm{G}} / \mathrm{A}\) is generated by the loops \(p(f)\) , for \(f\in \mathrm{S}\) , and the group \((\varpi_{\mathrm{G}} / \mathrm{A})_{p(a)}\) is generated by the elements of the form \(p(\theta_{\mathrm{G}}(f))\) , for \(f\in \mathrm{S} - \mathrm{Fl}(\mathrm{A})\) . The homomorphism \(p_a\circ \lambda\) is therefore surjective. Consequently, \(\lambda\) is surjective.

Remarks. --- 1) The homomorphism \(\mu_{a}\) is the inverse isomorphism of \(\lambda\) .

\begin{enumerate}
\def\labelenumi{\arabic{enumi})}
\setcounter{enumi}{1}
\tightlist
\item
  There exists a unique homomorphism \(\lambda\colon\mathrm{F}(\mathrm{Fl}(\mathrm{G}))\to\pi_{1}(\mathrm{G},a)\) defined by relations (1) for all \(f\in\mathrm{Fl}(\mathrm{G})\) . It follows from Proposition 8 that the homomorphism \(\lambda\) is surjective and its kernel is the smallest normal subgroup of \(\mathrm{F}(\mathrm{Fl}(\mathrm{G}))\) containing the elements \(f\) , for \(f\in\mathrm{Fl}(\mathrm{A})\) , and the elements \(f\cdot\overline{f}\) , for \(f\in\mathrm{Fl}(\mathrm{G})\) .
\end{enumerate}

\section{HOMOTOPIES}

\subsection{Definition of homotopies}

DEFINITION 1. --- Let G be a groupoid, let H be a quiver, and let \(\varphi\) and \(\varphi'\) be morphisms of quivers from H to G. A homotopy connecting \(\varphi\) to \(\varphi'\) is a map h from the vertex set of H to the set of arrows of G satisfying the following properties:

\begin{enumerate}
\def\labelenumi{(\roman{enumi})}
\item
  For every vertex a of H, the arrow \(h(a)\) has as its origin \(\varphi(a)\) and as its target \(\varphi'(a)\) ;
\item
  For every arrow \(f\) of H, with origin \(a\) and target \(b\) , we have \(\varphi(f)h(b)=h(a)\varphi'(f)\) .
\end{enumerate}

We say that \(\varphi\) and \(\varphi'\) are homotopic if there exists a homotopy connecting \(\varphi\) to \(\varphi'\) .

Let G be a groupoid, let H be a quiver and let \(\varphi\) , \(\varphi'\) , \(\varphi''\) be morphisms of quivers from H to G. The map \(a \mapsto e_{\varphi(a)}\) is a homotopy connecting \(\varphi\) to \(\varphi\) . If \(h\) is a homotopy connecting \(\varphi\) to \(\varphi'\) , the map \(a \mapsto h(a)^{-1}\) is a homotopy connecting \(\varphi'\) to \(\varphi\) . Let \(h\) and \(h'\) be homotopies connecting \(\varphi\) to \(\varphi'\) and \(\varphi'\) to \(\varphi''\) respectively. For every vertex \(a\) of H, the arrows \(h(a)\) and \(h'(a)\) are composable. The map \(a \mapsto h(a)h'(a)\) is a homotopy connecting \(\varphi\) to \(\varphi''\) .

It follows that the relation " \(\varphi\) is homotopic to \(\varphi'\) " is an equivalence relation on the set of quiver morphisms from H to G.

Let G be a groupoid, let H be a quiver, and let \(\varphi\) and \(\varphi'\) be morphisms of quivers from H to G that are homotopic. By condition (i) of Definition 1, for every vertex \(a\) of H, the vertices \(\varphi(a)\) and \(\varphi'(a)\) belong to the same orbit of G.

Suppose moreover that H is a groupoid and let h be a homotopy connecting \(\varphi\) to \(\varphi'\) . The maps \(\operatorname{Orb}(\varphi)\) and \(\operatorname{Orb}(\varphi')\) deduced from \(\varphi\) and \(\varphi'\) by passage to orbits are therefore equal. For every vertex a of H and every arrow \(f \in H_{a}\) , we have \(\varphi(f) = h(a)\varphi'(f)h(a)^{-1} = \operatorname{Int}(h(a))(\varphi'(f))\) , by condition (ii) of Definition 1. In other words, the homomorphism \(\varphi_{a}\) is equal to \(\operatorname{Int}(h(a)) \circ \varphi_{a}'\) . In particular, if the homomorphism \(\varphi_{a}\) is injective (resp. bijective, resp. surjective), then so is the homomorphism \(\varphi_{a}'\) .

Remarks. --- 1) Let G, G' be groupoids, let H, H' be quivers, let \(u\colon\mathrm{H}'\to\mathrm{H}\) be a morphism of quivers and \(v\colon\mathrm{G}\to\mathrm{G}'\) be a morphism of groupoids. If morphisms of quivers \(\varphi,\varphi'\) from H to G are homotopic, then the morphisms of quivers \(v\circ\varphi\circ u\) and \(v\circ\varphi'\circ u\) from H' to G' are homotopic. More precisely, if \(h\) is a homotopy connecting \(\varphi\) to \(\varphi'\) , the map \(\mathrm{Fl}(v)\circ h\circ\mathrm{Som}(u)\) is a homotopy connecting \(v\circ\varphi\circ u\) to \(v\circ\varphi'\circ u\) .

\begin{enumerate}
\def\labelenumi{\arabic{enumi})}
\setcounter{enumi}{1}
\tightlist
\item
  Let G be a groupoid, let H be a quiver and let \(\varphi, \psi\) be morphisms of quivers from H to G. Denote by \(j\) the canonical morphism from H to Grp(H), and let \(\overline{\varphi}\) and \(\overline{\psi}\) be the morphisms of groupoids from Grp(H) to G such that \(\overline{\varphi} \circ j = \varphi\) and \(\overline{\psi} \circ j = \psi\) .
\end{enumerate}

Recall that \(\operatorname{Som}(\mathrm{H}) = \operatorname{Som}(\operatorname{Grp}(\mathrm{H}))\) .

A homotopy \(h\colon \operatorname{Som}(\mathrm{H})\to \operatorname{Fl}(\mathrm{G})\) connecting \(\varphi\) to \(\psi\) is a homotopy connecting \(\overline{\varphi}\) to \(\overline{\psi}\) .

\subsection{Homotopisms of groupoids}

In this section, we shall use the notation \(u \sim v\) to express that two morphisms of groupoids u and v are homotopic.

DEFINITION 2. --- Let G, G' be groupoids and let \(\varphi\) be a morphism from G to G'. An inverse up to homotopy of \(\varphi\) is a morphism of groupoids \(\psi\) from G' to G such that the morphisms \(\psi \circ \varphi\) and \(\varphi \circ \psi\) are respectively homotopic to \(\mathrm{Id}_{\mathrm{G}}\) and to \(\mathrm{Id}_{\mathrm{G}'}\) . One says that \(\varphi\) is a homotopism if there exists an inverse of \(\varphi\) up to homotopy.

An isomorphism of groupoids is a homotopism.

Let G and G' be groupoids. Let \(\varphi\) , \(\varphi'\) be morphisms of groupoids from G to G' that are homotopic. If \(\varphi\) is a homotopism, then so is \(\varphi'\) . Indeed, if \(\psi\) denotes a homotopy inverse of \(\varphi\) , we have \(\psi \circ \varphi' \sim \psi \circ \varphi \sim \mathrm{Id}_{\mathrm{G}}\) and \(\varphi' \circ \psi \sim \varphi \circ \psi \sim \mathrm{Id}_{\mathrm{G}'}\) , which proves that \(\psi\) is a homotopy inverse of \(\varphi'\) .

Let G, G', and G'' be groupoids and let \(\varphi\colon\mathrm{G}\to\mathrm{G}'\) , \(\varphi'\colon\mathrm{G}'\to\mathrm{G}''\) , \(\psi\colon\mathrm{G}'\to\mathrm{G}\) , \(\psi'\colon\mathrm{G}''\to\mathrm{G}'\) be morphisms of groupoids. Then, among the following conditions:

\begin{enumerate}
\def\labelenumi{(\roman{enumi})}
\item
  \(\psi\) is a homotopy inverse of \(\varphi\) ;
\item
  \(\psi'\) is a homotopy inverse of \(\varphi'\) ;
\item
  \(\psi \circ \psi'\) is a homotopy inverse of \(\varphi' \circ \varphi\) ;
\end{enumerate}

any two imply the third. Indeed, suppose first that (i) and (ii) are satisfied; we then have

\[
\psi \circ \psi^ {\prime} \circ \varphi^ {\prime} \circ \varphi \sim \psi \circ \mathrm{Id} _ {\mathrm{G} ^ {\prime}} \circ \varphi \sim \psi \circ \varphi \sim \mathrm{Id} _ {\mathrm{G}}
\]

and, similarly, \(\varphi' \circ \varphi \circ \psi \circ \psi' \sim \mathrm{Id}_{\mathrm{G}'}\) , whence (iii). If (i) and (iii) are satisfied,

\[
\varphi^ {\prime} \circ \psi^ {\prime} \sim \varphi^ {\prime} \circ \varphi \circ \psi \circ \psi^ {\prime} \sim \mathrm{Id} _ {\mathrm{G} ^ {\prime \prime}}
\]

and

\[
\psi^ {\prime} \circ \varphi^ {\prime} \sim (\varphi \circ \psi) \circ \psi^ {\prime} \circ \varphi^ {\prime} \circ (\varphi \circ \psi) \sim \varphi \circ \psi \sim \mathrm{Id} _ {\mathrm{G} ^ {\prime}},
\]

whence condition (ii). The proof that conditions (ii) and (iii) imply condition (i) is analogous.

In particular, if two of the morphisms \(\varphi\) , \(\varphi'\) , \(\varphi'\circ\varphi\) are homotopisms, the same holds for the third.

PROPOSITION 1. --- Let G, G' be groupoids, let \(\varphi\) be a morphism from G to G', and let A be a subset of the vertex set of G that meets every orbit of G. For \(\varphi\) to be a homotopism, it is necessary and sufficient that the following conditions be satisfied:

\begin{enumerate}
\def\labelenumi{(\roman{enumi})}
\item
  the map \(\mathrm{Orb}(\varphi)\) from \(\mathrm{Orb}(\mathrm{G})\) to \(\mathrm{Orb}(\mathrm{G}')\) , deduced from \(\varphi\) by passing to orbits, is bijective;
\item
  for all \(a \in \mathrm{A}\) , the homomorphism \(\varphi_a \colon \mathrm{G}_a \to \mathrm{G}_{\varphi(a)}'\) is bijective.
\end{enumerate}

Suppose first that \(\varphi\) is a homotopism and let \(\psi\) be an inverse of \(\varphi\) up to homotopy. Then,

\[
\operatorname{Orb} (\psi) \circ \operatorname{Orb} (\varphi) = \operatorname{Orb} (\psi \circ \varphi) = \operatorname{Orb} (\operatorname{Id} _ {\mathrm{G}}) = \operatorname{Id} _ {\operatorname{Orb} (\mathrm{G})},
\]

because two homotopic groupoid morphisms induce the same map by passing to orbits. Likewise, \(\mathrm{Orb}(\varphi) \circ \mathrm{Orb}(\psi) = \mathrm{Id}_{\mathrm{Orb}(G')}\) . The map \(\mathrm{Orb}(\varphi)\) is therefore bijective, whence assertion (i). We also have, for every vertex a of G,

\[
\psi_ {\varphi (a)} \circ \varphi_ {a} = (\psi \circ \varphi) _ {a};
\]

since \(\psi \circ \varphi\) is homotopic to \(\mathrm{Id}_{\mathrm{G}}\) , the homomorphism \((\psi \circ \varphi)_a\) is bijective (cf. p. 180), so that \(\varphi_a\) is injective and \(\psi_{\varphi(a)}\) is surjective. Exchanging the roles of \(\varphi\) and \(\psi\) , we see that \(\varphi_a\) is also surjective, whence condition (ii).

Now suppose that conditions (i) and (ii) are satisfied and let us prove that \(\varphi\) is a homotopism.

First treat the case where each orbit of G' is reduced to a point.

For each vertex \(b\) of \(G'\) , choose a vertex \(u(b)\) of \(G\) belonging to \(A\) and whose image under \(\varphi\) is \(b\) . This is possible because the map \(\operatorname{Orb}(\varphi)\) is surjective and \(A\) meets each orbit of \(G\) . Let \(f\) be an arrow of \(G'\) ; we have \(o(f) = t(f)\) by hypothesis; set \(b = o(f)\) and \(a = u(b)\) . By condition (ii), there exists a unique arrow \(v(f) \in G_a\) whose image under \(\varphi\) is \(f\) . The pair \(\psi = (u, v)\) is a morphism of groupoids from \(G'\) to \(G\) . Let us prove that \(\psi\) is an inverse of \(\varphi\) up to homotopy. We already have \(\varphi \circ \psi = \mathrm{Id}_{\mathrm{G}'}\) by construction of \(\psi\) .

Let \(x\) be a vertex of G. Set \(a = \psi(\varphi(x))\) ; this is an element of A such that \(\varphi(a) = \varphi(x)\) . Since \(\mathrm{Orb}(\varphi)\) is injective, \(a\) belongs to the orbit of \(x\) in G and there exists an arrow \(f\) in G connecting \(a\) to \(x\) . The arrow \(h(x) = \psi(\varphi(f))^{-1}f\) then connects \(a\) to \(x\) and one has \(\varphi(h(x)) = e_{\varphi(a)}\) .

Let us show that the map \(h \colon \operatorname{Som}(\mathbf{G}) \to \operatorname{Fl}(\mathbf{G})\) thus defined is a homotopy connecting \(\psi \circ \varphi\) to \(\operatorname{Id}_{\mathbf{G}}\) . Condition (i) of definition 1 is satisfied, by construction. Let \(f\) be an arrow of \(\mathbf{G}\) , \(x\) its origin and \(y\) its target. We have \(\varphi(x) = \varphi(y)\) ; set \(a = \psi(\varphi(x)) = \psi(\varphi(y))\) . The arrows \(h(x)fh(y)^{-1}\) and \(\psi \circ \varphi(f)\) belong to \(\mathrm{G}_a\) and both have image \(\varphi(f)\) in \(\mathrm{G}'_{\varphi(a)}\) . As the map \(\varphi_a\) is injective, we have \(h(x)fh(y)^{-1} = \psi \circ \varphi(f)\) . Condition (ii) of definition 1 is thus satisfied.

Now let us prove proposition 1 in the general case. Let X be a maximal oriented forest of G' (II, p. 157, prop. 1). Let G'' be the groupoid deduced from G' by contracting the arrows of X and let \(\varphi'\colon G' \to G''\) be the canonical morphism. The morphism \(\varphi'\) satisfies the conditions of prop. 1 (II, p. 170, remark 1 and p. 178, cor. 2), and so does the morphism \(\varphi'\circ\varphi\) .

Since the orbits of G'' are reduced to points, it follows from the particular case already proved that \(\varphi'\) and \(\varphi'\circ\varphi\) are homotopisms, so that \(\varphi\) is also one. This completes the proof of the proposition.

COROLLARY 1. --- Let G be a groupoid, let A be a set, and let \(f\colon A \to \operatorname{Som}(G)\) be a map. If the image of \(f\) meets each orbit of G, the canonical groupoid morphism \(\varphi\) from \(f^{*}G\) to G is a homotopism.

By definition of the inverse image groupoid (II, p. 166, example 4), A is the set of vertices of the groupoid \(f^{*}\mathrm{G}\) and one has \(\mathrm{Fl}_{a,b}(f^{*}\mathrm{G}) = \mathrm{Fl}_{f(a),f(b)}(\mathrm{G})\) for every pair \((a,b)\) of elements of A. Moreover, we have \(\operatorname{Som}(\varphi) = f\) and \(\operatorname{Fl}(\varphi)\) induces the identity map from \(\operatorname{Fl}_{a,b}(f^{*}\mathrm{G})\) to \(\operatorname{Fl}_{f(a),f(b)}(\mathrm{G})\) . Consequently, the map \(\operatorname{Orb}(\varphi)\) is bijective and the homomorphism \(\varphi_{a}\colon (f^{*}\mathrm{G})_{a} \to \operatorname{G}_{f(a)}\) is an isomorphism, for all \(a \in A\) . The hypotheses of Proposition 1 are therefore satisfied.

COROLLARY 2. --- Let G be a groupoid, let X be an oriented forest of G, and let G' be the groupoid obtained from G by contracting the arrows of X. The canonical morphism from G to G' is a homotopism.

This follows from Remark 1 of II, p. 170, and Corollary 2 of II, p. 178, which show that the hypotheses of Proposition 1 are satisfied.

\subsection{Cohomotoper}

Let H be a quiver, let G be a groupoid, and let \(\varphi\) and \(\psi\) be morphisms of quivers from H to G.

Denote by \(G_{1}\) the quiver defined in the following way: the vertices of \(G_{1}\) are those of G; the arrows of \(G_{1}\) are the elements of the disjoint union of the sets \(\mathrm{Fl(G)}\) and \(\mathrm{Som(H)}\) ; the source map of \(G_{1}\) coincides with that of G on \(\mathrm{Fl(G)}\) and with \(\mathrm{Som}(\varphi)\) on \(\mathrm{Som(H)}\) ; the target map of \(G_{1}\) coincides with that of G on \(\mathrm{Fl(G)}\) and with \(\mathrm{Som}(\psi)\) on \(\mathrm{Som(H)}\) . Denote by \(\alpha_{1}\) the quiver morphism from G to \(G_{1}\) defined by the identity map of \(\mathrm{Som(G)}\) and by the canonical injection of \(\mathrm{Fl(G)}\) into \(\mathrm{Fl(G_{1})}\) . Denote by \(h_{1}\) the canonical injection \(\mathrm{Som(H)} \to \mathrm{Fl(G_{1})}\) .

Consider the free groupoid \(\mathrm{Grp}(G_{1})\) constructed on \(G_{1}\) (II, p. 174, def. 9) and denote by \(\theta_{1}\) the canonical quiver morphism from \(G_{1}\) to \(\mathrm{Grp}(G_{1})\). Finally, denote by \(\mathrm{Coh}(\varphi,\psi)\) the groupoid obtained from \(\mathrm{Grp}(G_{1})\) by contracting the loops (at the origin of \(x\)).

\[
\alpha_ {1} (x) \alpha_ {1} (y) \alpha_ {1} (x y) ^ {- 1},\tag{1}
\]

for every pair \((x,y)\) of composable arrows of G, as well as loops (at \(\varphi(a)\))

\[
\alpha_ {1} (\varphi (f)) h _ {1} (b) \alpha_ {1} (\psi (f)) ^ {- 1} h _ {1} (a) ^ {- 1}\tag{2}
\]

for \(a\) , \(b\) in \(\operatorname{Som}(\mathrm{H})\) and \(f \in \operatorname{Fl}_{ab}(\mathrm{H})\) . Denote by \(\pi: \operatorname{Grp}(\mathrm{G}_1) \to \operatorname{Coh}(\varphi, \psi)\) the canonical morphism; set \(\alpha = \pi \circ \theta_1 \circ \alpha_1\) and \(h = \operatorname{Fl}(\pi \circ \theta_1) \circ h_1\) .

PROPOSITION 2. --- The groupoid Coh(φ,ψ) is generated by the subquiver whose set of vertices is Som(G) and whose set of arrows is the union of the images of the maps Fl(α) and h.

This subquiver being the image of the quiver \(G_{1}\) by the quiver morphism \(\pi \circ \theta_{1}\) , the proposition follows immediately from the construction of \(\mathrm{Coh}(\varphi, \psi)\) .

PROPOSITION 3. --- The quiver morphism \(\alpha\) is a groupoid morphism from G to \(\mathrm{Coh}(\varphi, \psi)\) and the map \(h\) is a homotopy connecting \(\alpha \circ \varphi\) to \(\alpha \circ \psi\) .

The triple \((\mathrm{Coh}(\varphi,\psi),\alpha,h)\) has the following universal property: if \(G'\) is a groupoid, \(\alpha'\) is a morphism of groupoids from G to \(G'\) and \(h'\colon\mathrm{Som}(\mathrm{H})\to\mathrm{Fl}(\mathrm{G}')\) is a homotopy connecting \(\alpha'\circ\varphi\) to \(\alpha'\circ\psi\) , there exists a unique morphism of groupoids \(\eta\colon\mathrm{Coh}(\varphi,\psi)\to\mathrm{G}'\) such that

\[
\alpha^ {\prime} = \eta \circ \alpha \qquad\text{and}\qquad h ^ {\prime} = \operatorname{Fl} (\eta) \circ h.\tag{3}
\]

Given the definition of the groupoid obtained by contracting arrows, the contraction of the loops (1) implies that \(\alpha\) is a morphism of groupoids, and the contraction of the loops (2) implies that \(h\) is a homotopy connecting \(\alpha \circ \varphi\) to \(\alpha \circ \psi\) .

Let \(G'\) , \(\alpha'\) , \(h'\) be as in the statement. Let \(\eta_{1}\) be the quiver morphism from \(G_{1}\) to \(G'\) such that \(\operatorname{Som}(\eta_{1})\) is equal to \(\operatorname{Som}(\alpha')\) and such that \(\operatorname{Fl}(\eta_{1})\) coincides with \(\operatorname{Fl}(\alpha')\) on \(\operatorname{Fl}(G)\) and with \(h'\) on \(\operatorname{Som}(H)\) . There exists a unique groupoid morphism \(\eta_{2}: \operatorname{Grp}(G_{1}) \to G'\) such that \(\eta_{1} = \eta_{2} \circ \theta_{1}\) . Since \(\alpha'\) is a morphism of groupoids and \(h'\) is a homotopy connecting \(\alpha' \circ \varphi\) to \(\alpha' \circ \psi\) , \(\eta_{2}\) defines by passage to the quotient a morphism of groupoids \(\eta\) from \(\operatorname{Coh}(\varphi, \psi)\) to \(G'\) (II, p. 170, prop. 3). This morphism satisfies relations (3) and is the only one (II, p. 185, prop. 2).

DEFINITION 3. --- The groupoid Coh(φ,ψ) is called the cohomotoper of the pair (φ,ψ). One says that α is the canonical morphism from G to Coh(φ,ψ) and that h is the canonical homotopy relating α∘φ to α∘ψ.

We call the frame of the pair \((\varphi,\psi)\) the quiver whose set of vertices is the set of orbits of G, whose set of arrows is the set of connected components of H, and whose source and target maps are induced by \(\varphi\) and \(\psi\) by passing to quotients.

PROPOSITION 4. --- The map Orb(α): Orb(G) → Orb(Coh(φ,ψ)) is surjective, and its fibers are the connected components of the frame of the pair (φ,ψ).

The morphism \(\alpha\) is the composite of the morphisms \(\alpha_{1}\colon\mathrm{G}\to\mathrm{G}_{1}\) , \(\theta_{1}\colon\mathrm{G}_{1}\to\mathrm{Grp}(\mathrm{G}_{1})\) and \(\pi\colon\mathrm{Grp}(\mathrm{G}_{1})\to\mathrm{Coh}(\varphi,\psi)\) . The map \(\theta_{1}\) induces a bijection from the set of connected components of the quiver \(\mathrm{G}_{1}\) to the set of orbits of \(\mathrm{Grp}(\mathrm{G}_{1})\) and the map \(\mathrm{Orb}(\pi)\colon\mathrm{Orb}(\mathrm{Grp}(\mathrm{G}_{1}))\to\mathrm{Orb}(\mathrm{Coh}(\varphi,\psi))\) is bijective (II, p. 170, remark 1). It therefore suffices to prove that the map from \(\mathrm{Orb}(\mathrm{G})\) to \(\pi_{0}(\mathrm{G}_{1})\) induced by \(\alpha_{1}\) is surjective and that its fibers are the connected components of the frame of the pair \((\varphi,\psi)\) . Surjectivity follows from the fact that the map \(\mathrm{Som}(\alpha_{1})\) is the identity map. The equivalence relation in \(\mathrm{Som}(\mathrm{G})\) given by “\(a\) and \(b\) are in the same connected component of \(\mathrm{G}_{1}\)” is generated by the relations “there exists an arrow of G connecting \(a\) to \(b\)” and “there exists a vertex \(h\) of H such that \(\varphi(h)=a\) and \(\psi(h)=b\)”. This equivalence relation is compatible with the map from \(\mathrm{Som}(\mathrm{G})\) to \(\mathrm{Orb}(\mathrm{G})\), and the relation induced on \(\mathrm{Orb}(\mathrm{G})\) is generated by the relation “there exists an orbit \(\eta\) of H such that \(\mathrm{Orb}(\varphi)(\eta)=\alpha\) and \(\mathrm{Orb}(\psi)(\eta)=\beta\)”. This is therefore the relation “\(\alpha\) and \(\beta\) are in the same connected component of the frame of the pair \((\alpha,\beta)\)”.

PROPOSITION 5. --- Let G' be a groupoid, let η, η' be morphisms of groupoids from Coh(φ,ψ) to G' and let k be a map from Som(G) to Fl(G'). For k to be a homotopy connecting η to η', it is necessary and sufficient that the following two conditions be satisfied:

\begin{enumerate}
\def\labelenumi{(\roman{enumi})}
\item
  The map \(k\) is a homotopy connecting \(\eta \circ \alpha\) to \(\eta' \circ \alpha\) ;
\item
  For every vertex a of H, we have
\end{enumerate}

\[
\eta (h (a)) k (\psi (a)) = k (\varphi (a)) \eta^ {\prime} (h (a)).
\]

By definition, for k to be a homotopy connecting \(\eta\) to \(\eta'\) , it is necessary and sufficient that the following two conditions be satisfied (recall that \(\operatorname{Som}(G) = \operatorname{Som}(\operatorname{Coh}(\varphi, \psi))\) ) :

\begin{enumerate}
\def\labelenumi{\alph{enumi})}
\item
  For every vertex \(x\) of \(\operatorname{Coh}(\varphi, \psi)\) , \(k(x)\) connects \(\eta(x)\) to \(\eta'(x)\) ;
\item
  For every pair \((x,y)\) of vertices of \(\mathrm{Coh}(\varphi,\psi)\) and every arrow \(f\in\mathrm{Fl}_{x,y}(\mathrm{Coh}(\varphi,\psi))\) , we have \(\eta(f)k(y)=k(x)\eta'(f)\) .
\end{enumerate}

By prop. 2, it suffices to verify condition \(b\) ) when \(f\) belongs to the image of \(\mathrm{Fl}(\alpha)\) or to that of \(h\) , so that \(b\) ) is equivalent to the conjunction of the two conditions \(c\) ) and \(d\) ) below:

\begin{enumerate}
\def\labelenumi{\alph{enumi})}
\setcounter{enumi}{2}
\item
  For \(x \in \operatorname{Som}(\mathrm{G})\) , \(y \in \operatorname{Som}(\mathrm{G})\) and \(g \in \operatorname{Fl}_{x,y}(\mathrm{G})\) , we have \(\eta(\alpha(g))k(y) = k(x)\eta'(\alpha(g))\) ;
\item
  For every \(a \in \operatorname{Som}(\mathrm{H})\) , we have \(\eta(h(a))k(\psi(a)) = k(\varphi(a))\eta'(h(a))\) . Condition (i) is equivalent to the conjunction of a) and c), and condition (ii) is condition d), whence the corollary.
\end{enumerate}

\subsection{Comparison of two cohomotopers}

Consider a diagram

\[
\begin{array}{c} \mathrm{H} \xrightarrow {\varphi} \mathrm{G} \\ u \Big \downarrow \qquad \qquad \qquad \qquad \qquad \qquad \qquad \qquad \qquad \Big \downarrow v \\ \mathrm{H} ^ {\prime} \xrightarrow {\varphi^ {\prime}} \mathrm{G} ^ {\prime} \\ \psi^ {\prime} \end{array}\tag{4}
\]

where H, H', G, G' are groupoids and u, v, \(\varphi\) , \(\psi\) , \(\varphi'\) , \(\psi'\) are morphisms of groupoids such that \(v \circ \varphi = \varphi' \circ u\) and \(v \circ \psi = \psi' \circ u\) .

Denote by \(\alpha\) the canonical morphism from G to the cohomotoper \(\mathrm{Coh}(\varphi, \psi)\) and \(h\) the canonical homotopy connecting \(\alpha \circ \varphi\) to \(\alpha \circ \psi\) ; define \(\alpha'\) and \(h'\) analogously. Then, \(\alpha' \circ v\) is a groupoid morphism from G to \(\mathrm{Coh}(\varphi', \psi')\) and \(h' \circ \mathrm{Som}(u)\) is a homotopy connecting \(\alpha' \circ \varphi' \circ u\) to \(\alpha' \circ \psi' \circ u\) , that is, \(\alpha' \circ v \circ \varphi\) to \(\alpha' \circ v \circ \psi\) . By the universal property of cohomotopers (II, p. 185, prop. 3), there exists a unique morphism of groupoids \(w\) from \(\mathrm{Coh}(\varphi, \psi)\) to \(\mathrm{Coh}(\varphi', \psi')\) such that \(w \circ \alpha = \alpha' \circ v\) and \(\mathrm{Fl}(w) \circ h = h' \circ \mathrm{Som}(u)\) . We have in particular extended diagram (4) to a diagram

\[
\begin{array}{c} \mathrm{H} \xrightarrow [ \psi ]{\varphi} \mathrm{G} \xrightarrow {\alpha} \mathrm{Coh} (\varphi , \psi) \\ u \Biggl \downarrow \qquad \qquad \qquad \Biggl \downarrow v \qquad \qquad \Biggl \downarrow w \\ \mathrm{H} ^ {\prime} \xrightarrow [ \psi^ {\prime} ]{\varphi^ {\prime}} \mathrm{G} ^ {\prime} \xrightarrow {\alpha^ {\prime}} \mathrm{Coh} (\varphi^ {\prime}, \psi^ {\prime}) \end{array}\tag{5}
\]

in which the second square is commutative.

THEOREM 1. --- Assume the following hypotheses:

\begin{enumerate}
\def\labelenumi{(\roman{enumi})}
\item
  the morphism of groupoids v is a homotopism;
\item
  the map \(\operatorname{Orb}(u)\colon\operatorname{Orb}(\mathrm{H})\to\operatorname{Orb}(\mathrm{H}')\) , deduced from u by passing to orbits, is bijective;
\item
  there exists in each orbit of H a point a such that the homomorphism \(u_a \colon \mathrm{H}_a \to \mathrm{H}_{u(a)}'\) is surjective.
\end{enumerate}

Then, the morphism of groupoids \(w\colon \operatorname{Coh}(\varphi,\psi)\to\operatorname{Coh}(\varphi',\psi')\) is a homotopism.

Let G'' be the groupoid deduced from G' by contracting the arrows of a maximal oriented forest. The canonical morphism \(v'\colon G' \to G''\) is a homotopism (II, p. 184, corollary 2 of prop. 1). The two diagrams

\[
\begin{array}{c c c} \mathrm{H} \xrightarrow [ \psi ]{\varphi} \mathrm{G} & & \mathrm{H} ^ {\prime} \xrightarrow [ \psi^ {\prime} ]{\varphi^ {\prime}} \mathrm{G} ^ {\prime} \\ u \Big \downarrow & \Big \downarrow_ {v ^ {\prime} \circ v} & \text {and} \quad \mathrm{Id} _ {\mathrm{H} ^ {\prime}} \Big \downarrow \\ \mathrm{H} ^ {\prime} \xrightarrow [ v ^ {\prime} \circ \psi^ {\prime} ]{v ^ {\prime} \circ \varphi^ {\prime}} \mathrm{G} ^ {\prime} & & \mathrm{H} ^ {\prime} \xrightarrow [ v ^ {\prime} \circ \psi^ {\prime} ]{v ^ {\prime} \circ \varphi^ {\prime}} \mathrm{G} ^ {\prime \prime} \end{array}
\]

give rise to morphisms of groupoids \(w_1' \colon \mathrm{Coh}(\varphi, \psi) \to \mathrm{Coh}(v' \circ \varphi', v' \circ \psi')\) and \(w_2' \colon \mathrm{Coh}(\varphi', \psi') \to \mathrm{Coh}(v' \circ \varphi', v' \circ \psi')\) ; we have \(w_1' = w_2' \circ w\) . It therefore suffices to prove that \(w_1'\) and \(w_2'\) are homotopisms. Since the maps \(\mathrm{Som}(v' \circ v) \colon \mathrm{Som}(\mathrm{G}) \to \mathrm{Som}(\mathrm{G}'' )\) and \(\mathrm{Som}(v') \colon \mathrm{Som}(\mathrm{G}') \to \mathrm{Som}(\mathrm{G}'' )\) are surjective, it consequently suffices to prove the theorem under the additional hypothesis that the map \(\mathrm{Som}(v)\) is surjective, a hypothesis that we shall make throughout the rest of the proof.

Let us prove successively the following assertions:

- The map Orb(w) is bijective;

- For every vertex \(a\) of G, the homomorphism \(w_{a}\) is surjective;

- For every vertex \(a\) of G, the homomorphism \(w_{a}\) is injective.

\begin{enumerate}
\def\labelenumi{\alph{enumi})}
\item
  By hypothesis, the map \(\operatorname{Orb}(u)\) is bijective; the same holds for the map \(\operatorname{Orb}(v)\) according to II, p. 182, prop. 1, since v is a homotopism. The quiver morphism from the frame of the pair \((\varphi,\psi)\) to that of the pair \((\varphi',\psi')\) defined by the maps \(\operatorname{Orb}(u)\) and \(\operatorname{Orb}(v)\) is therefore an isomorphism. In particular, the map deduced from it by passing to connected components is bijective. Prop. 4 of II, p. 185 then implies that the map \(\operatorname{Orb}(w)\) is bijective.
\item
Let \(f'\) be an arrow of \(G'\) , denote by \(a'\) its origin and \(b'\) its target. Since the map \(\operatorname{Som}(v)\) is surjective, there exist vertices \(a\) and \(b\) in G such that \(a' = v(a)\) and \(b' = v(b)\) . Since the morphism v is a homotopism, there exists an arrow \(f\) of G connecting \(a\) to \(b\) , as well as an element \(g \in G_{a}\) such that \(v(g) = f'v(f)^{-1}\) (II, p. 182, prop. 1), whence \(f' = v(gf)\) . This proves that the map \(\mathrm{Fl}(v)\) is surjective.
\end{enumerate}

Let us then prove that the map \(\mathrm{Fl}(w)\) is surjective. Its image contains that of \(\mathrm{Fl}(\alpha')\) , because one has \(\alpha'\circ v=w\circ\alpha\) and the map \(\mathrm{Fl}(v)\) is surjective. Let b be a vertex of \(H'\) ; let a be a vertex of H such that b and \(u(a)\) are in the same orbit of \(H'\) and let f be an arrow of \(H'\) connecting \(u(a)\) to b. Then,

\[
h ^ {\prime} (u (a)) \cdot (\alpha^ {\prime} \circ \psi^ {\prime}) (f) = (\alpha^ {\prime} \circ \varphi^ {\prime}) (f) \cdot h ^ {\prime} (b),
\]

since \(h'\) is a homotopy connecting \(\alpha' \circ \varphi'\) to \(\alpha' \circ \psi'\) . The arrow \(h'(u(a)) = w(h(a))\) belongs to the image of \(\mathrm{Fl}(w)\) , as do the two arrows \(\alpha'(\psi'(f))\) and \(\alpha'(\varphi'(f))\) by what precedes. It follows that the arrow \(h'(b)\) belongs to the image of \(\mathrm{Fl}(w)\) , which proves that the image of \(\mathrm{Fl}(w)\) contains that of \(h'\) . By prop. 2 of II, p. 185, the map \(\mathrm{Fl}(w)\) is surjective.

Let \(g' \in \mathrm{Coh}(\varphi', \psi')_{u(a)}\) . Let g be an arrow of \(\mathrm{Coh}(\varphi, \psi)\) such that \(w(g) = g'\) . Denote by x and y the source and target of g; we have \(u(x) = u(y) = u(a)\) . Let \(g_1\) (resp. \(g_2\) ) be an arrow of G connecting a to x (resp. a to y) whose image under v is \(e_{u(a)}\) . Then, \(\alpha(g_1)g\alpha(g_2)^{-1}\) is an element of \(\mathrm{Coh}(\varphi, \psi)_a\) whose image under \(w_a\) is \(g'\) . Consequently, for every vertex a of G, the homomorphism \(w_a\) is surjective.

\begin{enumerate}
\def\labelenumi{\alph{enumi})}
\setcounter{enumi}{2}
\tightlist
\item
  Let us prove that, for every vertex \(a\) of G, the homomorphism \(w_{a}\) is injective. Considering successively the diagrams
\end{enumerate}

\[
\begin{array}{c c c} \mathrm{H} \xrightarrow {\varphi} \mathrm{G} & & \mathrm{H} \xrightarrow {v \circ \varphi} \mathrm{G} ^ {\prime} \\ \mathrm{Id} _ {\mathrm{H}} \Biggl \downarrow \quad \psi & \Biggl \downarrow v & u \Biggl \downarrow \quad \Biggl \downarrow \mathrm{Id} _ {\mathrm{G} ^ {\prime}} \\ \mathrm{H} \xrightarrow {v \circ \varphi} \mathrm{G} ^ {\prime} & & \mathrm{H} ^ {\prime} \xrightarrow {\varphi^ {\prime}} \mathrm{G} ^ {\prime} \\ & v \circ \psi & \psi^ {\prime} \end{array}
\]

We are then reduced to treating the following two cases: 1) We have H' = H and \(u = \mathrm{Id}_{\mathrm{H}}\) ; 2) We have G' = G and \(v = \mathrm{Id}_{\mathrm{G}}\) .

\begin{enumerate}
\def\labelenumi{\arabic{enumi})}
\tightlist
\item
  Suppose that H' = H and \(u = \mathrm{Id}_{\mathrm{H}}\) .
\end{enumerate}

Consider a morphism of groupoids \(v^{\prime}: G^{\prime} \to G\) which is inverse to v up to homotopy and a homotopy \(k: \operatorname{Som}(G) \to \operatorname{Fl}(G)\) connecting \(v^{\prime} \circ v\) to \(\mathrm{Id}_{\mathrm{G}}\) . By Remark 1 of II, p. 181, the maps \(\alpha \circ k \circ \varphi\) and \(\alpha \circ k \circ \psi\) are homotopies connecting respectively \(\alpha \circ v^{\prime} \circ \varphi^{\prime}\) to \(\alpha \circ \varphi\) and \(\alpha\circ v^{\prime}\circ\psi^{\prime}\) to \(\alpha\circ\psi\) . Consequently (cf. II, p. 180), the map

\[
\begin{array}{c} h _ {1} \colon \operatorname{Som} (\mathrm{H}) \to \operatorname{Fl} (\operatorname{Coh} (\varphi , \psi)), \\ x \mapsto (\alpha \circ k \circ \varphi) (x) \cdot h (x) \cdot ((\alpha \circ k \circ \psi) (x)) ^ {- 1} \end{array}
\]

is a homotopy connecting \(\alpha \circ v' \circ \varphi'\) to \(\alpha \circ v' \circ \psi'\) . By the universal property of cohomotopers (II, p. 185, prop. 3), there exists a unique morphism of groupoids \(w'\colon \mathrm{Coh}(\varphi', \psi') \to \mathrm{Coh}(\varphi, \psi)\) such that \(\alpha \circ v' = w' \circ \alpha'\) and \(h_1 = \mathrm{Fl}(w') \circ h'\) .

\[
\begin{array}{c} \mathrm{H} \xrightarrow [ \psi ]{\varphi} \mathrm{G} \xrightarrow {\alpha} \mathrm{Coh} (\varphi , \psi) \\ \mathrm{Id} _ {\mathrm{H}} \Biggl \downarrow \qquad \qquad \qquad \Biggl \downarrow v \qquad \qquad \Biggl \downarrow w \\ \mathrm{H} \xrightarrow [ \psi^ {\prime} ]{\varphi^ {\prime}} \mathrm{G} ^ {\prime} \xrightarrow {\alpha^ {\prime}} \mathrm{Coh} (\varphi^ {\prime}, \psi^ {\prime}) \\ \mathrm{Id} _ {\mathrm{H}} \Biggl \downarrow \qquad \qquad \qquad \Biggl \downarrow v ^ {\prime} \qquad \qquad \Biggl \downarrow w ^ {\prime} \\ \mathrm{H} \xrightarrow [ \psi ]{\varphi} \mathrm{G} \xrightarrow {\alpha} \mathrm{Coh} (\varphi , \psi)  . \end{array}
\]

In particular, we have

\[
\alpha \circ v ^ {\prime} \circ v = w ^ {\prime} \circ \alpha^ {\prime} \circ v = w ^ {\prime} \circ w \circ \alpha .
\]

Since k is a homotopy connecting \(v'\circ v\) to \(\mathrm{Id}_{\mathrm{G}}\) , \(\alpha\circ k\) is a homotopy connecting \(w'\circ w\circ\alpha\) to \(\alpha\) . Since \(\mathrm{Fl}(w'\circ w)\circ h=\mathrm{Fl}(w')\circ h'=h_{1}\) , we have, for every vertex x of H,

\[
\operatorname{Fl} \left(w ^ {\prime} \circ w\right) \circ h (x) \cdot (\alpha \circ k \circ \psi) (x) = h _ {1} (x) \cdot (\alpha \circ k \circ \psi) (x) = (\alpha \circ k \circ \varphi) (x) \cdot h (x),
\]

by definition of \(h_1\) . By prop. 5 of II, p. 186, applied to the morphisms of groupoids \(w' \circ w\) and \(\mathrm{Id}_{\mathrm{Coh}(\varphi, \psi)}\) , the map \(\alpha \circ k\) is a homotopy connecting \(w' \circ w\) to \(\mathrm{Id}_{\mathrm{G}}\) . In particular, \(w' \circ w\) is a homotopism.

For every vertex \(a\) of G, the group homomorphism \((w' \circ w)_a\) is therefore bijective (II, p. 182, prop. 1). It follows that the homomorphism \(w_a\) is injective, whence the result in case A).

\begin{enumerate}
\def\labelenumi{\arabic{enumi})}
\setcounter{enumi}{1}
\tightlist
\item
  Suppose we have G' = G and \(v = \mathrm{Id}_{\mathrm{G}}\) .
\end{enumerate}

Let \(x\) be a vertex of H'. The map \(\operatorname{Orb}(u)\) being surjective, there exists a vertex \(a\) of H and an arrow \(f\) of H' connecting \(u(a)\) to \(x\) . The arrows \(\alpha(\varphi'(f))^{-1}\) , \(h(a)\) and \(\alpha(\psi'(f))\) connect respectively \(\varphi'(x)\) to \(\varphi'(u(a)) = \varphi(a)\) , \(\varphi(a)\) to \(\psi(a)\) and \(\psi'(u(a)) = \psi(a)\) to \(\psi'(x)\) ; they are therefore composable in \(\operatorname{Coh}(\varphi, \psi)\) . Set

\[
h _ {2} (x) = \alpha (\varphi^ {\prime} (f)) ^ {- 1} \cdot h (a) \cdot \alpha (\psi^ {\prime} (f)).
\]

Let us verify that the arrow \(h_{2}(x)\) thus defined does not depend on the chosen elements a and f. Let \(a'\) be a vertex of H and let \(f'\) be an arrow of \(H'\) connecting \(u(a')\) to x. Since the map \(\operatorname{Orb}(u)\) is injective, the vertices a and \(a'\) of H belong to the same orbit and there exists an arrow \(c \in \operatorname{Fl}(\operatorname{H})\) connecting a to \(a'\) . Then, \(u(c)f'f^{-1}\) is a loop at \(u(a)\) in \(H'\) . By hypothesis (iii), there exists a vertex b of H such that the homomorphism \(u_{b}\) is surjective and an arrow \(c'\) of H from b to a; then, \(\operatorname{Int}(u(c'))(u(c)f'f^{-1})\) is a loop at \(u(b)\) in H, and it is therefore the image under \(u_{b}\) of a loop \(c''\) at b. Consequently, the arrow \(g = (c'c)^{-1}c''c'\) of H connects the vertex \(a'\) to the vertex a and satisfies \(f' = u(g)f\) . One then has

\[
\begin{array}{r l} & {\alpha (\varphi^ {\prime} (f ^ {\prime})) ^ {- 1} h (a ^ {\prime}) \alpha (\psi^ {\prime} (f ^ {\prime}))} \\ & {\qquad = \alpha (\varphi^ {\prime} (f)) ^ {- 1} \alpha (\varphi^ {\prime} (u (g))) ^ {- 1} h (a ^ {\prime}) \alpha (\psi^ {\prime} (u (g))) \alpha (\psi^ {\prime} (f))} \\ & {\qquad = \alpha (\varphi^ {\prime} (f)) ^ {- 1} \alpha (\varphi (g)) ^ {- 1} h (a ^ {\prime}) \alpha (\psi (g)) \alpha (\psi^ {\prime} (f))} \\ & {\qquad = \alpha (\varphi^ {\prime} (f)) ^ {- 1} \cdot h (a) \cdot \alpha (\psi^ {\prime} (f))} \end{array}
\]

since h is a homotopy connecting \(\alpha \circ \varphi\) to \(\alpha \circ \psi\) . This proves the announced independence.

By construction, we have \(h_{2}(u(x)) = h(x)\) for all \(x \in \operatorname{Som}(\mathrm{H})\) . We have thus defined a map \(h_{2}\) from \(\operatorname{Som}(\mathrm{H}')\) to \(\operatorname{Fl}(\operatorname{Coh}(\varphi, \psi))\) . Let c be an arrow of \(H'\) , denote its source by x and its target by y. Let a be a vertex of H and let f be an arrow of \(H'\) connecting \(u(a)\) to x. Then fc is an arrow of \(H'\) which connects \(u(a)\) to y. By definition of \(h_{2}\) , we thus have

\[
\begin{array}{r l} & h _ {2} (x) \alpha (\psi^ {\prime} (c)) = \alpha (\varphi^ {\prime} (f)) ^ {- 1} h (a) \alpha (\psi^ {\prime} (f)) \alpha (\psi^ {\prime} (c)) \\ & \qquad = \alpha (\varphi^ {\prime} (c)) \alpha (\varphi^ {\prime} (f c)) ^ {- 1} h (a) \alpha (\psi^ {\prime} (f c)) \\ & \qquad = \alpha (\varphi^ {\prime} (c)) h _ {2} (y). \end{array}
\]

This proves that \(h_{2}\) is a homotopy connecting \(\alpha \circ \varphi'\) to \(\alpha \circ \psi'\) .

By the universal property of cohomotopers, there exists a unique groupoid morphism \(w'\colon \mathrm{Coh}(\varphi', \psi') \to \mathrm{Coh}(\varphi, \psi)\) such that \(w' \circ \alpha' = \alpha\) and \(h_2 = \mathrm{Fl}(w') \circ h'\) . We have \(w' \circ w \circ \alpha = w' \circ \alpha' = \alpha\) and \(\mathrm{Fl}(w' \circ w) \circ h = \mathrm{Fl}(w') \circ h' \circ \mathrm{Som}(u) = h_2 \circ u = h\) by definition of \(h_2\) . By the universal property of cohomotopers, this implies that \(w' \circ w = \mathrm{Id}_{\mathrm{Coh}(\varphi, \psi)}\) . In particular, for every \(a \in \mathrm{Som}(G)\) , the homomorphism \(w_a\) is injective.

It then follows from prop. 1 of II, p. 182 that the morphism w is a homotopism, whence the theorem.
